% Part: sets-functions-relations
% Chapter: size-of-sets-complete

\documentclass[../../../include/open-logic-chapter]{subfiles}

\begin{document}

\olchapter{sfr}{siz}{La taille des ensembles}

\begin{editorial}
Ce chapitre traite des énumérations, de la dénombrabilité et de la
non-dénombrabilité. Plusieurs sections existent en deux versions :
une présentation plus élémentaire, où les énumérations sont des
listes, ou des surjections depuis $\PosInt$, et une présentation
plus abstraite, où les énumérations sont des bijections avec
$\Nat$ ou un segment initial.\footnote{Précision éditoriale :
l'original abrège ici la seconde définition en ne mentionnant que
$\Nat$ ; les segments initiaux sont nécessaires pour les ensembles
finis et figurent dans la définition détaillée.}
\end{editorial}

\olimport{introduction}

\olimport{enumerability}

\olimport{zig-zag}

\olimport{pairing}

\olimport{pairing-alt}

\olimport{non-enumerability}

\olimport{reduction}

\olimport{equinumerous-sets}

\olimport{comparing-size}

\olimport{schroder-bernstein}

\begin{editorial}
La \olref[sfr][siz][enm-alt]{sec}, la
\olref[sfr][siz][nen-alt]{sec} et la \olref[sfr][siz][red-alt]{sec}
qui suivent sont des variantes de la
\olref[sfr][siz][enm]{sec}, de la \olref[sfr][siz][nen]{sec} et
de la \olref[sfr][siz][red]{sec}, rédigées par Tim Button pour son
ouvrage \emph{Open Set Theory}. Un peu plus avancées, elles
utilisent une autre définition des énumérations, plus adaptée
au contexte de la théorie des ensembles : une bijection avec
$\Nat$ ou avec un segment initial, au lieu d'une présentation
par liste ou par image d'une surjection depuis~$\PosInt$.
\end{editorial}

\olimport{enumerability-alt}
\olimport{non-enumerability-alt}
\olimport{reduction-alt}

\OLEndChapterHook

\end{document}
